\begin{abstract}
An Internet access point serves applications that differ in kind: video calls and games need a modest rate in every second, while downloads and mail need volume and can wait.
The operator also sells free, premium and pay-as-you-go access on the same link, and knows neither who will arrive nor how much capacity or demand the next second brings.
Protecting real-time traffic when capacity is short is the natural policy, but the class of a flow is in effect declared by the device, and any rule that rewards the declaration ``real-time'' invites every application to make it.
We propose \repay{}, an online allocation rule that grants real-time priority on credit.
Each delay-tolerant device that is passed over accumulates a deficit that raises its own priority and the price of further priority for the real-time flows of its tier, so that the volume is paid back when capacity returns.
\repay{} is a drift-plus-penalty method in which tier-level virtual queues enforce long-run business constraints and per-device deficit queues, created at arrival, enforce fairness between service classes; each per-slot problem is solved exactly by bisection.
We prove that every deficit is bounded by a constant on every sample path, for arbitrary arrivals, departures and capacities, so that no delay-tolerant device envies the real-time treatment by more than a fixed volume over any window of its stay; that for price-taking devices a static misreport of the class gains at most this constant in volume, and that what a switching strategy collects while it poses as real-time is at most this constant per switch; and that, for the fluid allocation and capacity that is independent across slots given the population, the average utility is at most $O(1/V)$ below that of a benchmark of policies that satisfy the business constraints and are envy-free between classes in expectation in every population state.
In simulations of the downlink of a public hotspot with about $75$ devices, on synthetic traces, \repay{} serves $90\%$ of real-time slots against $82\%$ for proportional fairness, at a delay-tolerant throughput within $2\%$ of it; strict priority serves all of them and the same scheduler without deficits $92\%$, but both reward misreporting. Strict priority with a well-chosen cap on the real-time share does better than \repay{} on both averages there and also blunts misreporting, without any guarantee to individual devices. In a trace-driven scenario built from public Wi-Fi session logs, where only one session in ten is real-time, \repay{} still gains $7$ points over proportional fairness, but the priority rules serve $12$ to $15$ points more real-time slots at a small average cost to the others and a class envy some sixty to over a hundred times larger: the case for \repay{} is its bound for every device, whatever the traffic mix, and it has a price.
Over $180$ devices on six traces, the median device that falsely declares video gains $170\%$ in volume under strict priority, $154\%$ under weighted fair queueing and $62\%$ under utility-maximising scheduling without deficits; under \repay{} the median change is $0.2\%$; about half of the devices that stay at least a minute still gain, by small amounts and none by more than a third of the bound, and most devices that stay less than a minute gain, by a median of $10\%$. In the trace-driven scenario the median change under \repay{} is $-0.2\%$ against $+24\%$ to $+66\%$ under the other class-aware rules, but $10$ of $120$ devices exceed the bound, one of them in a thin premium tier by a factor of ten, which the price-taking analysis does not exclude.
\end{abstract}

\noindent\textbf{Keywords:} online resource allocation; network utility maximisation; drift-plus-penalty; fairness; envy-freeness; incentive compatibility; quality of service; access networks.
